\documentclass[a4paper]{article}

\usepackage{graphicx}
\usepackage{amsmath,amssymb}
\usepackage{minted}
\usepackage{multirow}
\usepackage{float}
\usepackage{algorithm}
\usepackage{algpseudocode}
\usepackage{todonotes}
\usepackage{forest}
\usepackage{xstring}
\usepackage{xcolor}
\usepackage[most]{tcolorbox}
\usepackage{xparse}
\usepackage{natbib}
\usepackage{adjustbox}

\usetikzlibrary{graphs,graphdrawing}
\usegdlibrary{layered}

% for external graphviz
\input{/home/donkarlo/Dropbox/repo/nd_language_ai_project/src/nd_language_ai/formal/computer/programming/tex/latex/code_clips/external_graphviz.tex}

% for tags
\input{/home/donkarlo/Dropbox/repo/nd_language_ai_project/src/nd_language_ai/formal/computer/programming/tex/latex/part/control_sequence/kind/semtag/semtag.tex}

% for links
\input{/home/donkarlo/Dropbox/repo/nd_language_ai_project/src/nd_language_ai/formal/computer/programming/tex/latex/code_clips/seehere.tex}

% for nested lists and sections
\input{/home/donkarlo/Dropbox/repo/nd_language_ai_project/src/nd_language_ai/formal/computer/programming/tex/latex/code_clips/deep_structure.tex}

\title{Tight Fragment Preview and Extended Zoom}
\date{2026-09-08}
\author{Mohammad Rahmani}

\begin{document}
   \maketitle

   LaTeX fragment previews now use the \texttt{preview} package in tight-page mode so the generated PDF page follows the natural content bounds instead of clipping wide content at the A4 page boundary. This keeps included graphics and resolved child content outside scaling boxes and leaves complete standalone LaTeX documents on their original page geometry.

   The PDF preview maximum custom zoom factor was increased from 500\% to 2000\% so large diagrams can be inspected at substantially higher magnification.

   Regression coverage verifies that fragment previews use tight-page rendering without \texttt{adjustbox}, while complete standalone documents remain outside the tight-page fragment wrapper.

   \bibliographystyle{plainnat}
   \bibliography{/home/donkarlo/Dropbox/repo/nd_shared_research/bibliography/bibliography.bib}
\end{document}
